% =====================================================================
\chapter{Momentum and Nesterov}
\label{ch:momentum}
% =====================================================================

\newcommand{\tempnoisy}{\addplot[inkgreen,thin] coordinates {(1,14.29) (3,11.77) (5,10.48) (7,8.08) (9,10.05) (11,11.92) (13,11.24) (15,10.15) (17,12.75) (19,13.67) (21,12.33) (23,13.32) (25,12.32) (27,15.88) (29,12.56) (31,11.70) (33,14.71) (35,13.62) (37,9.13) (39,12.30) (41,14.46) (43,13.49) (45,14.43) (47,16.82) (49,15.07) (51,14.62) (53,16.07) (55,14.98) (57,15.88) (59,16.80) (61,13.27) (63,14.54) (65,13.01) (67,17.59) (69,13.19) (71,17.35) (73,19.56) (75,17.77) (77,19.63) (79,16.97) (81,20.44) (83,19.13) (85,18.66) (87,17.00) (89,20.65) (91,17.67) (93,19.16) (95,16.97) (97,19.12) (99,20.51) (101,18.95) (103,19.22) (105,18.50) (107,19.75) (109,20.43) (111,22.62) (113,17.10) (115,19.53) (117,17.00) (119,22.49) (121,19.34) (123,21.11) (125,24.44) (127,22.51) (129,24.22) (131,23.70) (133,21.32) (135,24.14) (137,21.34) (139,22.71) (141,20.54) (143,26.06) (145,21.31) (147,20.15) (149,23.72) (151,22.97) (153,22.72) (155,21.83) (157,21.37) (159,21.37) (161,24.10) (163,22.54) (165,23.86) (167,19.52) (169,22.12) (171,24.26) (173,24.61) (175,25.54) (177,22.76) (179,22.11) (181,24.63) (183,22.96) (185,23.26) (187,22.05) (189,22.95) (191,21.90) (193,21.91) (195,21.77) (197,23.11) (199,23.11) (201,25.41) (203,23.07) (205,20.14) (207,19.85) (209,22.57) (211,22.43) (213,23.04) (215,22.11) (217,22.11) (219,20.91) (221,18.74) (223,20.15) (225,24.01) (227,22.11) (229,20.11) (231,20.94) (233,21.14) (235,21.63) (237,22.33) (239,19.39) (241,20.97) (243,23.76) (245,19.23) (247,20.62) (249,20.29) (251,21.26) (253,23.94) (255,21.34) (257,20.99) (259,22.68) (261,19.75) (263,18.37) (265,20.95) (267,21.15) (269,19.09) (271,19.83) (273,19.87) (275,18.92) (277,19.52) (279,17.85) (281,17.80) (283,14.90) (285,19.56) (287,19.86) (289,16.41) (291,14.65) (293,19.84) (295,19.69) (297,18.14) (299,12.77) (301,16.48) (303,13.83) (305,14.76) (307,14.51) (309,14.66) (311,17.67) (313,13.23) (315,14.81) (317,15.09) (319,13.77) (321,13.54) (323,12.04) (325,13.15) (327,12.21) (329,13.39) (331,14.87) (333,14.45) (335,11.94) (337,11.62) (339,13.97) (341,14.91) (343,10.48) (345,10.49) (347,14.38) (349,11.78) (351,13.11) (353,12.60) (355,15.23) (357,9.91) (359,10.62) (361,11.36) (363,10.01) (365,6.92)};}
\newcommand{\tempewc}{\addplot[inkred,very thick] coordinates {(1,14.29) (3,10.93) (5,10.67) (7,10.21) (9,10.30) (11,11.39) (13,11.33) (15,11.11) (17,11.33) (19,11.63) (21,11.74) (23,12.22) (25,12.18) (27,12.70) (29,12.66) (31,12.75) (33,12.95) (35,13.14) (37,12.87) (39,13.04) (41,13.02) (43,13.27) (45,13.32) (47,13.77) (49,14.09) (51,14.37) (53,14.47) (55,14.67) (57,14.75) (59,14.67) (61,14.70) (63,14.81) (65,14.86) (67,15.14) (69,15.24) (71,15.51) (73,15.96) (75,16.07) (77,16.38) (79,16.54) (81,16.78) (83,17.22) (85,17.60) (87,17.50) (89,17.78) (91,17.63) (93,17.83) (95,17.92) (97,17.89) (99,18.45) (101,18.61) (103,18.80) (105,18.97) (107,19.14) (109,19.43) (111,20.18) (113,20.10) (115,20.02) (117,19.85) (119,20.36) (121,20.11) (123,20.17) (125,20.78) (127,20.96) (129,21.34) (131,21.70) (133,21.49) (135,21.66) (137,21.73) (139,21.78) (141,21.57) (143,22.13) (145,22.02) (147,21.68) (149,22.02) (151,22.15) (153,22.25) (155,22.02) (157,22.02) (159,21.94) (161,22.18) (163,22.21) (165,22.54) (167,22.51) (169,22.63) (171,22.84) (173,23.18) (175,23.42) (177,23.27) (179,23.25) (181,23.67) (183,23.85) (185,23.66) (187,23.55) (189,23.49) (191,23.52) (193,23.46) (195,23.10) (197,22.91) (199,23.08) (201,23.38) (203,23.33) (205,22.92) (207,22.70) (209,22.77) (211,22.88) (213,22.75) (215,22.57) (217,22.51) (219,22.31) (221,21.91) (223,21.73) (225,22.15) (227,21.84) (229,21.65) (231,21.66) (233,21.35) (235,21.37) (237,21.29) (239,21.21) (241,21.09) (243,21.27) (245,21.06) (247,20.88) (249,21.13) (251,21.10) (253,21.32) (255,21.14) (257,21.05) (259,20.95) (261,20.70) (263,20.24) (265,20.32) (267,20.42) (269,20.27) (271,20.23) (273,20.05) (275,19.99) (277,19.76) (279,19.71) (281,19.36) (283,18.96) (285,18.96) (287,18.81) (289,18.52) (291,18.04) (293,18.05) (295,17.97) (297,17.86) (299,17.23) (301,16.99) (303,16.42) (305,16.26) (307,15.79) (309,15.78) (311,16.06) (313,15.70) (315,15.35) (317,15.40) (319,15.24) (321,15.04) (323,14.62) (325,14.19) (327,13.80) (329,13.84) (331,14.33) (333,14.16) (335,13.89) (337,13.74) (339,13.24) (341,13.26) (343,12.81) (345,12.76) (347,12.89) (349,12.72) (351,12.69) (353,12.48) (355,12.75) (357,12.30) (359,12.15) (361,12.03) (363,11.91) (365,11.32)};}
\newcommand{\tempew}{\addplot[inkblue,very thick] coordinates {(1,1.43) (3,2.96) (5,4.37) (7,5.33) (9,6.31) (11,7.82) (13,8.45) (15,8.82) (17,9.44) (19,10.06) (21,10.45) (23,11.13) (25,11.31) (27,11.97) (29,12.07) (31,12.26) (33,12.55) (35,12.82) (37,12.61) (39,12.83) (41,12.85) (43,13.13) (45,13.21) (47,13.67) (49,14.01) (51,14.31) (53,14.42) (55,14.63) (57,14.71) (59,14.64) (61,14.68) (63,14.79) (65,14.85) (67,15.13) (69,15.23) (71,15.50) (73,15.95) (75,16.06) (77,16.38) (79,16.53) (81,16.78) (83,17.22) (85,17.60) (87,17.50) (89,17.78) (91,17.62) (93,17.83) (95,17.92) (97,17.89) (99,18.45) (101,18.61) (103,18.80) (105,18.97) (107,19.14) (109,19.43) (111,20.18) (113,20.10) (115,20.02) (117,19.85) (119,20.36) (121,20.11) (123,20.17) (125,20.78) (127,20.96) (129,21.34) (131,21.70) (133,21.49) (135,21.66) (137,21.73) (139,21.78) (141,21.57) (143,22.13) (145,22.02) (147,21.68) (149,22.02) (151,22.15) (153,22.25) (155,22.02) (157,22.02) (159,21.94) (161,22.18) (163,22.21) (165,22.54) (167,22.51) (169,22.63) (171,22.84) (173,23.18) (175,23.42) (177,23.27) (179,23.25) (181,23.67) (183,23.85) (185,23.66) (187,23.55) (189,23.49) (191,23.52) (193,23.46) (195,23.10) (197,22.91) (199,23.08) (201,23.38) (203,23.33) (205,22.92) (207,22.70) (209,22.77) (211,22.88) (213,22.75) (215,22.57) (217,22.51) (219,22.31) (221,21.91) (223,21.73) (225,22.15) (227,21.84) (229,21.65) (231,21.66) (233,21.35) (235,21.37) (237,21.29) (239,21.21) (241,21.09) (243,21.27) (245,21.06) (247,20.88) (249,21.13) (251,21.10) (253,21.32) (255,21.14) (257,21.05) (259,20.95) (261,20.70) (263,20.24) (265,20.32) (267,20.42) (269,20.27) (271,20.23) (273,20.05) (275,19.99) (277,19.76) (279,19.71) (281,19.36) (283,18.96) (285,18.96) (287,18.81) (289,18.52) (291,18.04) (293,18.05) (295,17.97) (297,17.86) (299,17.23) (301,16.99) (303,16.42) (305,16.26) (307,15.79) (309,15.78) (311,16.06) (313,15.70) (315,15.35) (317,15.40) (319,15.24) (321,15.04) (323,14.62) (325,14.19) (327,13.80) (329,13.84) (331,14.33) (333,14.16) (335,13.89) (337,13.74) (339,13.24) (341,13.26) (343,12.81) (345,12.76) (347,12.89) (349,12.72) (351,12.69) (353,12.48) (355,12.75) (357,12.30) (359,12.15) (361,12.03) (363,11.91) (365,11.32)};}

\pgfplotsset{tempaxis/.style={width=10cm,height=5cm,axis lines=left,xmin=0,xmax=380,ymin=5,ymax=28,
  xtick={0,170,365},ytick=\empty,xlabel={$t$ (days)},ylabel={$\theta$ ($^\circ$C)},
  xlabel style={font=\small,at={(1,-0.12)},anchor=north east},ylabel style={font=\small},tick label style={font=\scriptsize}}}

\section{Exponentially Weighted Average}

\lettrine[lines=2,findent=2pt,nindent=0pt]{\color{inkred}L}{et's} say we want to represent the temperature in Casablanca for the whole year:

\begin{center}
\begin{tikzpicture}\begin{axis}[tempaxis]\tempnoisy\end{axis}\end{tikzpicture}
\end{center}

When we \hlg{interpolate (green curve)} the data points, we get a noisy pattern. But even with this noisy pattern, we can globally know the global pattern of the curve: it's increasing, reaching a maximum, then decreasing.

So, to avoid noises and keep only the pattern, we should remark something:

\Rem[Where do noises come from?] When we try to interpolate, we link data points, for example:
\begin{center}
\begin{tikzpicture}
  \draw[inkgreen,thick] (0,0) -- (0.4,0.9) -- (0.8,0.35) -- (1.2,1.1) -- (1.5,0.2);
  \foreach \p in {(0,0),(0.4,0.9),(0.8,0.35),(1.2,1.1),(1.5,0.2)} \fill \p circle (1.5pt);
  \node[font=\small\itshape,anchor=west] (n) at (2.4,0.6) {This provokes noises: brutal deviations that come at once.};
  \draw[-Stealth] (n.west) -- (1.25,1.0);
\end{tikzpicture}
\end{center}

So, if we find a solution to avoid considering brutal data deviations, by giving much importance to the past evolution\ldots\ let me introduce you:

\Def{Def: Exponentially Weighted Average}
EWA is defined as an iterative weighted average:
\begin{keyformula}
\[
v_t=\beta\,v_{t-1}+(1-\beta)\,\theta_t
\]
\end{keyformula}
where $(v_t)$ is the representative sequence, and $\theta_t$ the function that we want to make smooth, while $\beta$ is a number $\in\,]0,1[$ that manages the degree of smoothness. But the main thing is that \textbf{$\beta$ controls how many steps we want to memorize from the past}, so as to avoid brutal deviation problems.

\begin{center}
\begin{tikzpicture}\begin{axis}[tempaxis]\tempnoisy\tempewc\end{axis}\end{tikzpicture}
\end{center}
The \hlr{red curve} represents the EWA effect for $\beta=0.9$.

\Rem As a rule of thumb, very used in ML and DL:
\[
N=\frac{1}{1-\beta}\qquad\text{with $N$: the number of remembered steps.}
\]
For example, $\beta=0.9$: $N=\frac{1}{1-0.9}=10$ steps.

And we can assume that:
\[
\boxed{\beta=1-\frac1N}\implies\text{if we want to remember $N$ steps, this is $\beta$.}
\]
For example, we want to remember 100 steps: $\beta=1-\frac{1}{100}=0.99$.

\hl{But let us see something very interesting:}
\begin{align*}
v_{t-1}&=\beta v_{t-2}+(1-\beta)\theta_{t-1}\\
v_t&=\beta v_{t-1}+(1-\beta)\theta_t\\
v_{t+1}&=\beta v_t+(1-\beta)\theta_{t+1}\\
\implies v_{t+1}&=(1-\beta)\theta_{t+1}\ \underset{\times\beta}{\longrightarrow}\ (1-\beta)\beta\,\theta_t\ \underset{\times\beta}{\longrightarrow}\ (1-\beta)\beta^2\,\theta_{t-1}+\cdots
\end{align*}
(the three terms are summed). So we multiply by $\beta$ at each iteration.

\Q{But when is it forgotten?}
\[
\beta=1-\frac1N\implies\beta^N=\left(1-\frac1N\right)^N,\qquad
\lim_{N\to+\infty}\beta^N=\lim_{N\to+\infty}\left(1-\frac1N\right)^N=\frac1e
\]
So if we choose $\beta$ depending on $N$, $\beta=1-\frac1N$, then the contribution after $N$ steps is $\frac1e\approx37\%$.

\begin{goldbox}
That's where the ``exponential'' comes from in EWA.
\end{goldbox}

\subsection*{\color{inkblue}Remark: Bias Correction}

The red curve shown earlier isn't just an EWA, but an EWA \emph{with bias correction}:

\begin{center}
\begin{tikzpicture}\begin{axis}[tempaxis]\tempnoisy\tempew\tempewc\end{axis}\end{tikzpicture}
\end{center}

The \hl{blue curve} shows what happens when we don't do bias correction, and the reason is that, theoretically and also in practice, we initialize $v_0=0$. But the problem is:
\[
v_1=0+(1-\beta)\,\theta_1
\]
In our example, where $\theta_t$ is the temperature and $\beta=0.9$, let's say $\theta_1=11^\circ$C. Then $v_1=0.1\times11=1.1^\circ\text{C}\ll11^\circ$C, so we start far away from the starting point.

To correct that, we change the EWA formula:
\begin{keyformula}[This is Bias Correction]
\[
\hat v_t=\frac{v_t}{1-\beta^t}=\frac{\beta v_{t-1}+(1-\beta)\theta_t}{1-\beta^t}
\]
\end{keyformula}
With $t=1$: $\hat v_1=\frac{1.1}{1-0.9}=11^\circ$C, exactly $\theta_1$. And since $\beta^t\to0$, the correction disappears by itself after a few steps.

People don't mind fixing that in practice. Even libraries like PyTorch or TensorFlow don't add bias correction for momentum, since data scientists consider it as a ``warm-up''.\ednote{In fact PyTorch avoids the problem differently: the momentum buffer is initialised with the first gradient ($v_1=g_1$) instead of $0$.} But for optimizers that use multiple moving averages (like Adam), bias correction is \textbf{mandatory}.\ednote{The notes cite ``RMSProp/Adam''. Adam applies bias correction (it's in the algorithm, see Chapter~\ref{ch:adam}), but the classical RMSProp, and its PyTorch/TensorFlow implementations, do not.}

\section{Momentum}

\Def[Persist in directions that keep reinforcing themselves.]{Def: Momentum}

As you may have remarked, the squiggly line represented before looks the same as the SGD noisy pattern. So we apply EWA on the $\nabla f_i(x_k)$ (in practice $\nabla f_T(x_k)$ for mini-batch GD).

So the algorithm is as follows:
\begin{keyformula}
\[
\left\{\begin{aligned}
v_t&=\beta\,v_{t-1}+(1-\beta)\,\nabla f_t(x_t)\\
x_{t+1}&=x_t-\gamma\,v_t
\end{aligned}\right.
\]
\end{keyformula}

\Rem Some AI engineers, and some frameworks like PyTorch, use this form of GD with momentum:
\[
\left\{\begin{aligned}
v_t&=\beta\,v_{t-1}+\nabla f_t(x_t)\\
x_{t+1}&=x_t-\gamma\,v_t
\end{aligned}\right.
\]
And this is the reason why. Let's take our classical form, applied on GD ($g_t(x)=\nabla f_t(x)$):
\begin{align*}
v_t&=\beta v_{t-1}+(1-\beta)g_t(x)\\
\text{with:}\quad v_{t-1}&=\beta v_{t-2}+(1-\beta)g_{t-1}(x)\\
\text{so:}\quad v_t&=(1-\beta)\bigl(g_t(x)+\beta\,g_{t-1}(x)+\beta^2 g_{t-2}(x)+\cdots\bigr)
\end{align*}
So $g_t(x)=\nabla f_t(x)$ is scaled by:\ednote{The notes sum from $i=1$; the geometric sum $\sum\beta^i=\frac{1}{1-\beta}$ starts at $i=0$ (the current gradient has weight $\beta^0=1$). Same fix in the next formula.}
\[
(1-\beta)\sum_{i=0}^{+\infty}\beta^i=(1-\beta)\times\frac{1}{1-\beta}=1\qquad\text{(geometric series, converges because $\beta<1$)}
\]
\begin{itemize}
  \item So the gradient is always (approximately) scaled by \underline{1}.
\end{itemize}

Let's take the other formula, $v_t=\beta v_{t-1}+g_t(x)$. That leads to $\nabla f(x)$ being scaled by:
\[
\sum_{i=0}^{+\infty}\beta^i=\frac{1}{1-\beta}
\]
This is huge! Imagine if $\beta=0.9$: $\frac{1}{1-0.9}=10$, so the gradient is scaled by $10\times$, much faster! (Of course, we need to tune the learning rate.)

That's why engineers prefer this formula to the other one.

\begin{notebox}[Important remark]
As I said, in PyTorch we don't use the form with $1-\beta$, but that's the default case. The truth is that, for PyTorch, the formula is as follows:
\[
\left\{\begin{aligned}
v_t&=\beta\,v_{t-1}+(1-d)\,\nabla f_t(x_t)\\
x_{t+1}&=x_t-\gamma\,v_t
\end{aligned}\right.
\]
with $d$ the ``\textbf{dampening}'' parameter; by default, $d=0$.
\end{notebox}

\Q{You said before that noises could be helpful; why use an algorithm that smoothens noises?}

Well, good question! And the idea is, as I said, that the main consequence of noises is avoiding local minima and saddle point problems, by following another path rather than taking the straightforward one.

Well, SGD with momentum does follow a different path from the straightforward one, but rather than losing time switching directions at each iteration, we follow a more stable, smooth line that follows the global evolution of SGD.

\section{Nesterov}

\Def[See before you leap.]{Def: Nesterov}

Nesterov SGD is a smart idea found by Yurii Nesterov. The idea is: from the momentum formula,
\[
v_t=\beta v_{t-1}+\nabla f_t(x_t),\qquad\text{then:}\qquad x_{t+1}=x_t-\gamma v_t,
\]
let's see what happens:

\begin{center}
\begin{tikzpicture}
  \foreach \r in {0.6,1.25} \draw[very thick] (0.3,0) ellipse ({2.6*\r} and {\r});
  \draw[very thick] (0.3,0) ellipse (4.6 and 1.95);
  \fill[inkred] (0.9,0) circle (2.5pt);
  \coordinate (xt) at (-2.6,-0.2);
  \coordinate (mid) at (-1.6,1.05);
  \coordinate (xn) at (-0.65,0.55);
  \draw[inkblue,very thick,-Stealth] (xt) -- (mid);
  \draw[inkpurple,very thick,-Stealth] (mid) -- (xn);
  \draw[inkgreen,very thick,-Stealth] (xt) -- (xn);
  \node[inkblue,font=\small,anchor=east] at (xt) {$x_t$};
  \node[inkblue,font=\small,anchor=south west] at (xn) {$x_{t+1}$};
  \node[font=\small] at (0.3,-2.5) {{\color{inkgreen}$\bm{\wr\wr}$} $v_t$ \qquad {\color{inkblue}$\bm{\wr\wr}$} $\gamma\beta v_{t-1}$ \qquad {\color{inkpurple}$\bm{\wr\wr}$} $\gamma\nabla f_t(x_t)$};
\end{tikzpicture}
\end{center}

So, what we can conclude from the figure is that $\beta v_{t-1}$ already gives us an idea of the region where we would ``leap''. So why not take it into consideration, and adjust the gradient at that point?

So here comes the Nesterov solution:
\begin{keyformula}
\[
\begin{aligned}
\bar x_t&=x_t-\gamma\beta\,v_{t-1}&&\text{\color{inkgold}``See''}\\
v_t&=\beta v_{t-1}+\nabla f_t(\bar x_t)&&\text{\color{inkgold}``Before''}\\
x_{t+1}&=x_t-\gamma\,v_t&&\text{\color{inkgold}``You leap''}
\end{aligned}
\]
\end{keyformula}
